\VerdierPhysicalPageAnchor{0060}%
\VerdierSemanticAnchor{ch1-1-5-14-continuation}
Posons alors :

\VerdierSemanticAnchor{ch1-1-5-14-eq3}
\begin{equation}
  \begin{aligned}
    T(0)&=\operatorname{id}_{\mathcal C},\\
    T(n)&=T^n, && n>0,\\
    T(n)&=(T^{-1})^{-n}, && n<0.
  \end{aligned}
  \tag{1.5.14.3}
\end{equation}

\VerdierSemanticAnchor{ch1-1-5-14-cocycle-data}
Puis, pour tout couple $(n,n')$ d'entiers, définissons des
isomorphismes :
\[
  c(n',n)\colon T(n+n')\longrightarrow T(n)T(n')
\]
par :

\VerdierSemanticAnchor{ch1-1-5-14-eq4}
\begin{equation}
  \begin{aligned}
    n,n'&\geq 0,
      &c(n',n)&=\text{identité},\\
    n,n'&\leq 0,
      &c(n',n)&=\text{identité},\\
    n'&\leq 0, n\geq 0, n+n'\leq 0,
      &c(n',n)&=u^{(n)}\mathbin{*}T(n+n'),\\
    n'&\leq 0, n\geq 0, n+n'\geq 0,
      &c(n',n)&=T(n+n')\mathbin{*}u^{(-n')},\\
    n'&\geq 0, n\leq 0, n+n'\leq 0,
      &c(n',n)&=T(n+n')\mathbin{*}v^{(n')},\\
    n'&\geq 0, n\leq 0, n+n'\geq 0,
      &c(n',n)&=v^{(-n)}\mathbin{*}T(n+n').
  \end{aligned}
  \tag{1.5.14.4}
\end{equation}

\VerdierSemanticAnchor{ch1-1-5-14-conclusion}
Pour s'assurer qu'on a bien ainsi défini une $\mathbb Z$-catégorie, il
reste à vérifier qu'on a bien les identités du type (1.2.1.3). Cette
vérification s'appuie sur les identités classiques que vérifient les
morphismes $u^{(n)}$ et $v^{(n)}$. Nous laissons au lecteur le soin de
l'entreprendre.

\VerdierSemanticAnchor{ch1-1-6}
\subsection{Cocycle de Koszul. Règle des signes}

\VerdierSemanticAnchor{ch1-1-6-1}
\noindent\textbf{1.6.1.}\quad
Soit $e_1,\ldots,e_n$ une base de $\mathbb Z^n$. Désignons par
$\mathbb Z^n_{i,j}$, $i,j\in[1,n]$, $i\neq j$, le sous-groupe de
$\mathbb Z^n$ engendré par les éléments $e_i$ et $e_j$. Les
homomorphismes de restriction définissent un homomorphisme :

\VerdierSemanticAnchor{ch1-1-6-1-eq1}
\begin{equation}
  H^2(\mathbb Z^n,\mathbb Z/2\mathbb Z)
  \longrightarrow
  \bigoplus_{\substack{i,j\in[1,n]\\ i<j}}
    H^2(\mathbb Z^n_{i,j},\mathbb Z/2\mathbb Z).
  \tag{1.6.1.1}
\end{equation}

On a des isomorphismes (canoniques) :

\VerdierSemanticAnchor{ch1-1-6-1-eq2}
\begin{equation}
  H^2(\mathbb Z^n_{i,j},\mathbb Z/2\mathbb Z)
  \simeq \mathbb Z/2\mathbb Z.
  \tag{1.6.1.2}
\end{equation}

\VerdierSemanticAnchor{ch1-1-6-1-conclusion}
L'homomorphisme (1.6.1.1) est un isomorphisme.
